\documentclass[12pt]{article}

\usepackage[a4paper, margin=1in]{geometry}
\usepackage{amsmath,amssymb,amsthm}
\usepackage{mathtools}
\usepackage{mathrsfs}
\usepackage{enumitem}
\usepackage{xcolor}
\usepackage[hidelinks]{hyperref}
\usepackage{fontspec}

\defaultfontfeatures{Ligatures=TeX}
\IfFontExistsTF{Times New Roman}
  {\setmainfont{Times New Roman}}
  {\setmainfont{TeX Gyre Termes}}
\IfFontExistsTF{Menlo}
  {\setmonofont{Menlo}}
  {\setmonofont{Latin Modern Mono}}

\newtheorem{theorem}{Theorem}[section]
\newtheorem{proposition}[theorem]{Proposition}
\newtheorem{lemma}[theorem]{Lemma}
\newtheorem{corollary}[theorem]{Corollary}
\theoremstyle{definition}
\newtheorem{definition}[theorem]{Definition}
\theoremstyle{remark}
\newtheorem{remark}[theorem]{Remark}

\newcommand{\tightlist}{%
  \setlength{\itemsep}{0pt}\setlength{\parskip}{0pt}}

\title{%
  \textbf{Finite Self-Models and the Open Explanatory Gap:}\\
  \textbf{A Conditional Arithmetic Analysis}%
}

\author{
  Sidong Liu, PhD \\
  \small iBioStratix Ltd \\
  \small \texttt{sidongliu@hotmail.com}
}

\date{July 2026}

\begin{document}
\emergencystretch=2em
\raggedbottom
\hbadness=5000
\vbadness=5000

\maketitle

\begin{abstract}
Paper~G offers a programme-level diagnosis of possible self-reference
limits, while Paper~II retains a leading small-ball stress-tensor
extraction with an \(O(R^{d+3})\) remainder. Neither result supplies a
fixed-point no-go theorem, a physical Fisher spectrum, or a
self-description tower. This paper tests a proposed quantitative bridge
by defining a local Araki-relative-entropy reconstruction loss for a
stipulated finite statistical bottleneck.

The former local minimax drift-floor theorem is withdrawn. A finite
Fisher ceiling does not by itself imply a positive reconstruction loss:
on a short modular orbit, an admissible two-outcome channel can preserve
the orbit coordinate exactly and achieve zero loss. A future positive
risk theorem would therefore need a typed observation model that excludes
such lossless-orbit constructions and relates estimator risk to the
reconstruction loss.

The surviving formal result is a conditional arithmetic inversion. If a
stipulated positive sequence satisfies
\[
\log I_k=-\alpha k\log k+O(k)
\]
and the available budget grows at polynomial scale, then its maximal
admissible index obeys
\[
K(n)=O\!\left(\frac{\log n}{\log\log n}\right).
\]
This lemma proves neither a physical coefficient/Fisher spectrum nor a
self-model resource bound. The explanatory gap therefore remains open:
the paper identifies conditions required for a future quantitative
rebuild but does not establish a finite self-description horizon.
 \end{abstract}

\clearpage
\tableofcontents
\clearpage

\section{Introduction}\label{sec:1}

Paper~G\cite{PaperG} motivates the possibility that self-reference
limits an accessibility-based description. Its regress and collapse
scenarios are philosophical diagnostics, not a fixed-point no-go theorem
or an exhaustive dichotomy. Paper~III uses that discussion only as
motivation.

Paper~II\cite{PaperII} moved the discussion to a more concrete level.
In the higher-dimensional CHM ball framework, Paper~II retains the
leading stress-tensor extraction at order \(R^{d+1}\), with a remainder
\(O(R^{d+3})\). It does not construct independent coefficient
directions, a Fisher map, or a self-description tower from that
remainder.

The present paper defines a local Araki-relative-entropy reconstruction
loss for a stipulated finite statistical bottleneck, but withdraws the
claimed universal drift floor. Estimation risk in a specified noisy
experiment cannot be identified with reconstruction loss for a channel
that may preserve the orbit coordinate exactly. The surviving formal
result is instead a conditional arithmetic inversion for a stipulated
positive sequence and budget. The surrounding programme draws on:
\begin{enumerate}
\item the BKM/QFIM structure already used in Paper~I\cite{PaperI};
\item the self-referential Fisher and ego-rigidity formalism of
T-DOME~Paper~III\cite{TDOMEIII};
\item the leading CHM expansion and remainder order recorded in
Paper~II\cite{PaperII}.
\end{enumerate}

The arithmetic statement assumes, rather than derives, a positive
sequence \(I_k\) with
\[
\log I_k = -\alpha k\log k + O(k),
\]
and a polynomial-scale budget. Under those stipulations its maximal
index obeys
\[
K(n)=O\!\left(\frac{\log n}{\log\log n}\right).
\]
This proves neither a physical Fisher spectrum nor a self-model budget.

The trilogy-level interpretation also depends on Paper~I's supplied
shell interface \((Amb)\), its dynamical covariance axiom \((C)\), and
the scale-faithfulness axiom \((SFa)\). Modular flow supplies algebraic
time only; a phenomenal-time reading is conditional and interpretive.
T-DOME~Paper~III supplies an ego-rigidity prior model, not a theorem that
turns it into a reconstruction floor.

The paper is organized as follows. Section~\ref{sec:2} defines finite
observers, admissible self-models, and the self-description deficit.
Section~\ref{sec:3} records why the former drift-floor theorem is
withdrawn. Section~\ref{sec:4} proves the conditional arithmetic
inversion. Section~\ref{sec:5} retains only the exact CHM moment
coefficient and explains why it does not supply a Fisher spectrum.
Section~\ref{sec:6}
explains the trilogy-level significance of the result, and
Section~\ref{sec:7} states the remaining limitations and open problems.

\section{Finite Observers and the Self-Description Deficit}\label{sec:2}

\subsection{Finite Observers}\label{subsec:2.1}

Motivated by Paper~G's architectural discussion, we first specify what
counts as a finite statistical self-model. This definition does not
inherit a no-go theorem from Paper~G.

\begin{definition}[Finite observer]\label{def:2.1}
A finite observer of complexity \(n\) is a triple
\[
(\mathcal{A}_{\mathrm{obs}},\omega_{\mathrm{obs}},\mathfrak{M}_n),
\]
where:
\begin{enumerate}
\item \(\mathcal{A}_{\mathrm{obs}} \subset \mathcal{A}_c\) is the
observer's accessible subalgebra;
\item \(\omega_{\mathrm{obs}} = \omega|_{\mathcal{A}_{\mathrm{obs}}}\)
is the restricted state;
\item \(\mathfrak{M}_n\) is an \(n\)-dimensional statistical model used
internally to represent the observer's own accessible state.
\end{enumerate}
\end{definition}

The complexity parameter \(n\) measures effective statistical degrees of
freedom, not arbitrary description length. This is deliberate. Paper~III
is not about Kolmogorov complexity in full generality. It is about the
finite number of statistically usable directions available to an
observer trying to model itself.

\subsection{Modular Fisher Sensitivity}\label{subsec:2.1b}

The Tomita--Takesaki modular automorphism group
\(\sigma_t^\omega\) of a faithful normal state \(\omega\) on a von
Neumann algebra satisfies \(\omega\circ\sigma_t^\omega=\omega\) for
all \(t\). The orbit is therefore trivial at the level of states, and
the associated tangent vector vanishes. Since the modular orbit itself
carries no state-level Fisher information, we replace it with the
intrinsic accessible-unitary perturbation family as the operational
proxy for modular sensitivity.

\begin{definition}[Modular Fisher sensitivity]\label{def:2.1b}
Fix a self-adjoint element
\(h\in\mathcal{A}_{\mathrm{obs}}\) and define the perturbed family
\[
\omega_s^{(h)}
:=
\omega_{\mathrm{obs}}\circ\mathrm{Ad}(e^{ish}),
\qquad s\in(-\epsilon,\epsilon).
\]
The modular Fisher sensitivity of
\(\omega_{\mathrm{obs}}\) along \(h\) is
\[
\mathcal{I}_F(\omega_{\mathrm{obs}},h)
:=
\mathbf{Q}_{\omega_{\mathrm{obs}}}\!\left(
\frac{d\omega_s^{(h)}}{ds}\bigg|_{s=0},\;
\frac{d\omega_s^{(h)}}{ds}\bigg|_{s=0}
\right).
\]
The modular Fisher sensitivity of the observer is
\[
\mathcal{I}_F(\omega_{\mathrm{obs}})
:=
\sup\bigl\{
\mathcal{I}_F(\omega_{\mathrm{obs}},h)
:\;
h=h^*\in\mathcal{A}_{\mathrm{obs}},\;
\|h\|\le 1
\bigr\}.
\]
\end{definition}

The supremum is taken over unit-norm self-adjoint perturbations in the
observer's own accessible algebra. It is finite whenever the BKM form
is bounded on the unit ball of
\(\mathcal{A}_{\mathrm{obs}}\)-tangent directions, which holds in the
local regime used throughout this paper. The quantity
\(\mathcal{I}_F(\omega_{\mathrm{obs}})\) measures the observer's
maximal intrinsic sensitivity to unitary perturbations of its own
accessible state. It is strictly positive whenever
\(\omega_{\mathrm{obs}}\) is not a trace state on
\(\mathcal{A}_{\mathrm{obs}}\).

Throughout the paper, the term \emph{local modular drift family}
refers to the one-parameter family \(\{\omega_s^{(h)}\}\) for a fixed
perturbation \(h\). When the supremum in
Definition~\ref{def:2.1b} is attained, we take \(h\) to be a
maximizer; when it is not, we use a standard
\(\eta\)-optimization: for each \(\eta>0\), choose
\(h_\eta=h_\eta^*\in\mathcal{A}_{\mathrm{obs}}\) with
\(\|h_\eta\|\le 1\) and
\(\mathcal{I}_F(\omega_{\mathrm{obs}},h_\eta)\ge
\mathcal{I}_F(\omega_{\mathrm{obs}})-\eta\), carry out the proof for
\(h_\eta\), and let \(\eta\downarrow 0\). Despite the name, this
family is not the Tomita--Takesaki modular orbit (which is
state-preserving) but the accessible-unitary perturbation family that
serves as its operational proxy at the level of Fisher information.
This is the family along which the observer's self-description is
hardest to compress.

\subsection{Admissible Self-Models}\label{subsec:2.2}

The theorem chain does not require a literal embedding of the self-model
into a matrix algebra. That stronger route would force us to solve a
much harder operator-algebraic problem than the present paper needs.
What matters here is the existence of a finite statistical bottleneck.

\begin{definition}[Admissible self-model]\label{def:2.2}
An admissible \(n\)-dimensional self-model is a pair
\[
(\Pi_n,\Psi_n),
\]
where:
\begin{enumerate}
\item \(\Pi_n:\mathcal{S}(\mathcal{A}_{\mathrm{obs}})\to
\mathfrak{M}_n\) is the state-space map induced by a normal unital
completely positive map from a finite-dimensional algebra into
\(\mathcal{A}_{\mathrm{obs}}\), so that Petz monotonicity of relative
entropy applies;
\item \(\Psi_n:\mathfrak{M}_n\to
\mathcal{S}(\mathcal{A}_{\mathrm{obs}})\) is a recovery map;
\item \(\Pi_n\) factors through the observer's own accessible state
space, so no external ``God's-eye'' parameter is introduced;
\item near the reference state \(\omega_{\mathrm{obs}}\), the Fisher
matrix on \(\mathfrak{M}_n\) can be whitened:
\[
\mathbf{I}^{(n)}_\omega
=
\mathrm{diag}(\lambda_1,\dots,\lambda_n),
\qquad
\lambda_a \leq \mathcal{I}_F(\omega_{\mathrm{obs}}).
\]
\end{enumerate}
\end{definition}

The last condition is the substantive standing assumption of the
theorem, not a mere coordinate normalization. It says that the
compressed model has only \(n\) effective Fisher directions, and none of
them is sharper than the observer's own modular sensitivity scale. This
is the exact place where the theorem absorbs finite-dimensionality.

\subsection{Self-Description Deficit}\label{subsec:2.3}

The correct loss function is not a metric but a divergence. In a
Type~III setting there is no preferred trace-class density-matrix
picture, while Araki relative entropy is native to the operator
algebraic formulation \cite{Araki1976}.

\begin{definition}[Self-description deficit]\label{def:2.3}
Define the BKM ball of radius \(\delta\) around
\(\omega_{\mathrm{obs}}\) by
\[
\mathcal{B}_\delta(\omega_{\mathrm{obs}})
:=
\bigl\{
\omega\in\mathcal{S}(\mathcal{A}_{\mathrm{obs}})
:\;
\mathbf{Q}_{\omega_{\mathrm{obs}}}(\omega-\omega_{\mathrm{obs}},\,
\omega-\omega_{\mathrm{obs}})
<\delta^2
\bigr\}.
\]
Fix once and for all a \(\delta>0\) small enough that Araki relative
entropy admits its quadratic expansion on
\(\mathcal{B}_\delta(\omega_{\mathrm{obs}})\), and large enough that
the ego-rigidity prior \(\pi_{\mathrm{ego}}\) of
T-DOME~Paper~III\cite{TDOMEIII} can be supported on a segment of the
local modular drift family inside
\(\mathcal{B}_\delta(\omega_{\mathrm{obs}})\). The self-description
deficit of complexity \(n\) is
\[
\varepsilon(n)
:=
\inf_{(\Pi_n,\Psi_n)}
\sup_{\omega\in \mathcal{B}_\delta(\omega_{\mathrm{obs}})}
S\!\left(
\omega
\;\middle\|\;
\Psi_n\circ\Pi_n(\omega)
\right),
\]
where the infimum is taken over all admissible \(n\)-dimensional
self-models.
\end{definition}

Near a fixed state, Araki relative entropy has the familiar BKM/QFIM
quadratic expansion
\[
S(\omega\|\omega+\delta\omega)
=
\frac{1}{2}\mathbf{Q}_\omega(\delta\omega,\delta\omega)
+o(\|\delta\omega\|^2),
\]
which is the local information-geometric structure already used in
Paper~I\cite{PaperI} and in Petz's monotone-metric framework
\cite{Petz1996}. The minimax form is essential: it rules out trivial
constant recovery maps that memorize the reference state while failing
arbitrarily nearby. Thus \(\varepsilon(n)\) measures optimal local
uniform reconstruction loss in the intrinsic Type~III geometry of
states.

\section{The Fisher Ceiling for Admissible Self-Models}\label{sec:3}

\subsection{Petz Monotonicity and Finite Fisher Budget}\label{subsec:3.1}

The first ingredient is standard.

\begin{lemma}[Petz monotonicity]\label{lem:3.1}
For any normal coarse-graining \(\Pi_n\),
\[
S(\omega\|\sigma)\ge S(\Pi_n(\omega)\|\Pi_n(\sigma)).
\]
\end{lemma}

\begin{proof}
This is the monotonicity of relative entropy under coarse-graining
\cite{Petz1986}. The present paper uses it only in the local regime,
where it implies that compression cannot increase BKM/Fisher
information.
\end{proof}

Whitening then converts finite statistical dimension directly into two
distinct ceilings: one for the specific modular drift channel, and one
for the total trace budget across all retained directions.

\begin{proposition}[Directional and total Fisher ceilings]\label{prop:3.2}
Let \((\Pi_n,\Psi_n)\) be an admissible \(n\)-dimensional self-model.
Then:
\begin{enumerate}
\item along the local modular drift family through
\(\omega_{\mathrm{obs}}\), the Fisher information retained after
compression satisfies
\[
\mathcal{I}^{(n)}_{\mathrm{drift}}
\le
\mathcal{I}_F(\omega_{\mathrm{obs}});
\]
\item the total retained data Fisher information satisfies
\[
\mathcal{I}^{(n)}_{\mathrm{data}}
:=
\mathrm{tr}\,\mathbf{I}^{(n)}_\omega
\le
n\,\mathcal{I}_F(\omega_{\mathrm{obs}}).
\]
\end{enumerate}
\end{proposition}

\begin{proof}
\(\!\)(1) Apply Lemma~\ref{lem:3.1} to the one-parameter modular drift
family. In the local quadratic regime, monotonicity of relative entropy
implies monotonicity of the associated BKM/Fisher metric, so
compression cannot increase Fisher information for that specific scalar
parameter.

\(\!\)(2) In Fisher-whitened coordinates, the compressed Fisher matrix is
\[
\mathbf{I}^{(n)}_\omega
=
\mathrm{diag}(\lambda_1,\dots,\lambda_n),
\qquad
\lambda_a \le \mathcal{I}_F(\omega_{\mathrm{obs}}).
\]
Summing over the \(n\) effective directions gives
\[
\mathrm{tr}\,\mathbf{I}^{(n)}_\omega
=
\sum_{a=1}^n \lambda_a
\le
n\,\mathcal{I}_F(\omega_{\mathrm{obs}}).
\]
\end{proof}

\subsection{Withdrawal of the Drift-Floor Claim}\label{subsec:3.2}

The former local minimax drift-floor theorem and its corollary are
withdrawn. The argument applied a van Trees estimation-risk bound to a
compression/recovery loss without specifying a noisy observation model
that prevents the orbit coordinate from being retained.

The obstruction can be implemented by an internal finite-dimensional
channel. Let \(s\mapsto\omega_s\) be a nonconstant smooth local drift
orbit. Choose an effect \(0\le M\le\mathbf{1}\) in
\(\mathcal{A}_{\mathrm{obs}}\) with
\(\frac{d}{ds}\omega_s(M)|_{s=0}\ne0\). The normal unital completely
positive map
\[
\Phi:\mathbb{C}^2\longrightarrow\mathcal{A}_{\mathrm{obs}},
\qquad
\Phi(a,b)=aM+b(\mathbf{1}-M),
\]
induces the binary statistical coordinate
\[
\Pi_1(\omega_s)=\bigl(p(s),1-p(s)\bigr),
\qquad p(s):=\omega_s(M).
\]
By the inverse-function theorem, after restricting to a sufficiently
short orbit segment, \(p\) is injective. Define the recovery map on its
image, and extend it arbitrarily off that image, by
\[
\Psi_1\bigl(p(s),1-p(s)\bigr)=\omega_s.
\]
This code uses an observable in the accessible algebra rather than an
external parameter, and it factors through a finite-dimensional
commutative algebra as required by Definition~\ref{def:2.2}. Then
\[
S\!\left(\omega_s\middle\|\Psi_1\circ\Pi_1(\omega_s)\right)=0
\]
for every \(s\) on the encoded segment. A positive Bayes estimation
risk can arise only after a likelihood or sampling channel restricts
access to \(s\); it is not a lower bound on this lossless orbit
reconstruction. The previous proof therefore confused estimation risk
with reconstruction loss. Petz monotonicity and the stipulated Fisher
ceilings above remain bookkeeping facts within their stated model, but
they imply no nonzero reconstruction floor.

\section{A Conditional Arithmetic Inversion}\label{sec:4}

Paper~II\cite{PaperII} supplies a leading CHM coefficient and an
\(O(R^{d+3})\) remainder, not an infinite family of normalized Fisher
directions. We therefore separate the arithmetic statement from any
physical interpretation.

Let \(\{I_k\}_{k\ge1}\) be a \emph{stipulated} positive sequence and
let \(B(n)\ge1\) be a \emph{stipulated} budget. Define
\[
K(n):=\max\left\{K\ge1:
\sum_{j=1}^{K}I_j^{-1}\le B(n)\right\},
\]
with \(K(n)=0\) if the set is empty. No claim is made here that the
\(I_k\) are Fisher informations of physical coefficient directions or
that \(B(n)\) is a self-model resource.

\begin{lemma}[Conditional arithmetic inversion]\label{lem:arithmetic-asymptotic}
Suppose
\[
\log I_k=-\alpha k\log k+O(k),
\qquad \alpha>0,
\]
and \(\log B(n)=O(\log n)\). Then
\[
K(n)=O\!\left(\frac{\log n}{\log\log n}\right).
\]
\end{lemma}

\begin{proof}
If \(K=K(n)\ge1\), the final nonnegative summand satisfies
\(I_K^{-1}\le B(n)\). Hence
\[
\alpha K\log K+O(K)=-\log I_K\le\log B(n)=O(\log n).
\]
For sufficiently large \(K\), the \(K\log K\) term dominates the
linear remainder, so inversion gives
\(K=O(\log n/\log\log n)\). Bounded \(K\) obeys the same asymptotic
bound.
\end{proof}

This lemma is exactly the surviving claim. It establishes neither a
self-model budget nor a physical Fisher spectrum. In particular,
rescaling coefficient coordinates changes any candidate \(I_k\), so a
physical application requires a separately justified normalization,
likelihood, and Fisher map.

\section{Exact CHM Moment Without a Fisher-Spectrum Claim}
\label{sec:5}

The following exact integral is retained as an algebraic check. It does
not construct the coefficient directions or Fisher map absent from
Paper~II.

\subsection{Exact d=1 Signal Factor}\label{subsec:5.2}

In \(d=1\), the CHM modular Hamiltonian is
\[
K_B
=
\frac{\pi}{R}\int_{-R}^{R}(R^2-x^2)\,T_{00}(x)\,dx.
\]
Expanding \(T_{00}(x)\) around the ball center gives
\[
T_{00}(x)
=
\sum_{k\ge 0}\frac{x^{2k}}{(2k)!}\partial_x^{2k}T_{00}(0)+\cdots.
\]
Therefore the coefficient multiplying the \(2k\)-th derivative is
\[
c_k^{(d=1)}
=
\frac{\pi}{R}\cdot\frac{1}{(2k)!}
\int_{-R}^{R}(R^2-x^2)x^{2k}\,dx\cdot R^{-(2k+2)}.
\]
The radial moment is exact:
\[
\int_{-R}^{R}(R^2-x^2)x^{2k}\,dx
=
\frac{4R^{2k+3}}{(2k+1)(2k+3)}.
\]
Hence
\[
c_k^{(d=1)}
=
\frac{4\pi}{(2k+1)(2k+3)(2k)!},
\]
which shows exact factorial suppression:
\[
c_k^{(d=1)}
\sim
\frac{\mathrm{const}}{(2k)!\,k^2}.
\]

The factorial suppression of \(c_k^{(d=1)}\) is normalization-dependent
signal information. It does not by itself determine a fluctuation
kernel, a renormalized Fisher information \(I_k\), or a cumulative
inverse-information budget. Consequently this calculation does not
instantiate Lemma~\ref{lem:arithmetic-asymptotic}; it only shows what one
factor in a future, fully specified model would look like.

\section{Trilogy Implications}\label{sec:6}

\subsection{Paper G as Motivation Only}\label{subsec:6.1}

Paper~G\cite{PaperG} supplies philosophical motivation for studying
self-description, but no formal bridge maps its syntax-level scenarios
to the reconstruction loss defined here. The withdrawn drift-floor
argument does not quantify Paper~G, and this paper asserts no generic
fixed-point obstruction.

\subsection{Relocating the Gap in Paper II Terms}\label{subsec:6.2}

Paper~II\cite{PaperII} supplies a leading CHM stress-tensor term and an
\(O(R^{d+3})\) remainder. It does not supply normalized coefficient
directions or a Fisher map. Paper~III therefore does not reinterpret the
remainder as a physical control tower. Lemma~\ref{lem:arithmetic-asymptotic}
applies only after a positive sequence and budget are stipulated
independently.

\subsection{Trilogy-Level Interpretation}\label{subsec:6.3}

The trilogy-level reading is conditional on Paper~I's shell interface
\((Amb)\), dynamical axiom \((C)\), and scale-faithfulness assumption
\((SFa)\), as well as Paper~II's restricted CHM regime. These inputs
supply algebraic modular time and conditional geometric statements, not
a phenomenal-time identification. The surviving arithmetic lemma does
not locate or quantify an explanatory gap.

\section{Limitations and Open Problems}\label{sec:7}

The withdrawals identify four concrete tasks for any future rebuild.

\begin{enumerate}
\item \textbf{A typed observation model.} A future risk theorem must
specify a likelihood or sampling channel that prevents lossless orbit
coding, then distinguish estimator risk from state reconstruction loss.

\item \textbf{Normalized coefficient directions.} A physical use of
Lemma~\ref{lem:arithmetic-asymptotic} requires independently constructed
coefficient perturbations, a renormalized kernel, and a normalization
under which the Fisher sequence is invariantly defined.

\item \textbf{A Paper II bridge.} The remainder order alone is
insufficient. One must construct a map from small-ball coefficient data
to the statistical model before any resource comparison is meaningful.

\item \textbf{Ego-rigidity interface.} The T-DOME prior may enter a
specified Bayesian estimation problem, but it supplies no reconstruction
floor by itself. Its operational channel and boundary conditions must be
given before van Trees can be used.
\end{enumerate}

\section{Conclusion}\label{sec:8}

Paper~III tests a proposed quantitative route to the explanatory gap.
Paper~G supplies programme-level motivation, while Paper~II supplies a
leading stress-tensor extraction and a controlled remainder order; it
does not identify that remainder as a physical self-description carrier.
The present paper defines a finite-self-model reconstruction loss and
records why the former universal drift-floor inference fails.

A finite Fisher ceiling alone does not imply a nonzero local minimax
reconstruction loss. The lossless-orbit construction shows that a finite
channel can preserve the relevant modular-orbit coordinate exactly on a
short segment. A positive lower bound may still arise under an additional
typed observation model that excludes such encodings and connects
estimation risk to reconstruction loss, but a precise characterization
of those conditions remains open.

Separately, the conditional arithmetic lemma shows that a stipulated
positive sequence with
\(\log I_k=-\alpha k\log k+O(k)\) has maximal admissible index
\(K(n)=O(\log n/\log\log n)\) under a polynomial-scale budget. It does
not establish a physical Fisher spectrum, a self-model budget, or a
finite self-description horizon. The explanatory gap is therefore not
eliminated in general, but it is not quantitatively located by the
present results.

\appendix

\section{Trilogy Key Premises Summary}\label{app:premises}

The following list records the external inputs and their restricted
roles. Paper~III is not self-contained with respect to these trilogy
interfaces.

\begin{enumerate}
\item \textbf{Paper I algebraic interface} (Paper~I\cite{PaperI}).
The Araki relative entropy of two nearby faithful normal states on a
Type~III von Neumann algebra admits a local quadratic expansion
\(S(\omega\|\omega+\delta\omega)=\tfrac{1}{2}\mathbf{Q}_\omega
(\delta\omega,\delta\omega)+o(\|\delta\omega\|^2)\), where
\(\mathbf{Q}_\omega\) is the BKM/QFIM quadratic form. This is the
information-geometric structure used in Definitions~\ref{def:2.1b}
and~\ref{def:2.3} and throughout Section~\ref{sec:3}. The broader
trilogy interpretation also assumes Paper~I's shell interface \((Amb)\),
dynamical covariance \((C)\), and scale-faithfulness \((SFa)\). These
give conditional algebraic structure; modular flow is algebraic time,
not a derived phenomenal-time variable.

\item \textbf{Small-ball expansion with \(O(R^{d+3})\) remainder}
(Paper~II\cite{PaperII}).
In the CHM ball framework for conformal nets, the modular energy of a
small ball of radius \(R\) admits the expansion
\(\delta\langle K_{B,u}\rangle = \frac{2\pi V_d}{d+2}R^{d+1}
\delta\langle T_{\mu\nu}\rangle u^\mu u^\nu + O(R^{d+3})\).
The leading coefficient extracts the local stress tensor, while the
higher-order remainder is uncontrolled. Paper~II supplies neither an
infinite coefficient tower nor a Fisher map; Section~\ref{sec:5} uses
only a separate exact \(d=1\) moment calculation.

\item \textbf{Self-reference motivation} (Paper~G\cite{PaperG}).
Paper~G discusses regress and collapse scenarios but does not prove that
a self-referential fixed point generically fails to exist. Paper~III
uses it as motivation only and does not quantify it.

\item \textbf{Ego-rigidity prior}
(T-DOME~Paper~III\cite{TDOMEIII}).
The ego-rigidity prior \(\mathcal{I}_{\mathrm{ego}}\) is the Fisher
information of the observer's prior distribution over its own
self-model parameters:
\(\mathcal{I}_{\mathrm{ego}}:=\int(\partial_\sigma\log
\pi_{\mathrm{ego}}(\sigma))^2\pi_{\mathrm{ego}}(\sigma)\,d\sigma\).
It may enter a separately specified Bayesian likelihood model. It does
not lower-bound reconstruction loss without such an observation model,
and no theorem in this paper uses it to do so.
\end{enumerate}
 
\newpage
\addcontentsline{toc}{section}{References}

\begin{thebibliography}{99}

\bibitem{PaperI}
S.~Liu, \emph{The Inseparability of Geometry and Modular Flow: A Structural Inseparability Theorem and Its Implications for the Hard Problem}, Zenodo (2026), DOI: 10.5281/zenodo.19607387.

\bibitem{PaperII}
S.~Liu, \emph{Higher-Dimensional Consistency of Geometry and Modular Flow: Small-Ball Modular Data and Linearized Gravity}, Zenodo (2026), DOI: 10.5281/zenodo.19607543.

\bibitem{PaperG}
S.~Liu, \emph{Structural Limits of Unification: Accessibility, Incompleteness, and a Programme-Level Final-Cut Diagnosis}, Zenodo (2026), DOI: 10.5281/zenodo.18402907.

\bibitem{TDOMEIII}
S.~Liu, \emph{Fisher Information Geometry and Conditional Implementation Costs of Self-Referential Calibration}, Zenodo (2026), DOI: 10.5281/zenodo.18591771.

\bibitem{Araki1976}
H.~Araki,
\emph{Relative entropy of states of von Neumann algebras},
Publ.\ Res.\ Inst.\ Math.\ Sci.\ \textbf{11}, 809--833 (1976).

\bibitem{Petz1986}
D.~Petz,
\emph{Sufficient subalgebras and the relative entropy of states of a von Neumann algebra},
Commun.\ Math.\ Phys.\ \textbf{105}, 123--131 (1986).

\bibitem{Petz1996}
D.~Petz,
\emph{Monotone metrics on matrix spaces},
Linear Algebra Appl.\ \textbf{244}, 81--96 (1996).

\bibitem{vanTrees1968}
H.~L.~Van Trees,
\emph{Detection, Estimation, and Modulation Theory, Part~I},
Wiley, New York (1968).

\bibitem{CasiniHuertaMyers2011}
H.~Casini, M.~Huerta, and R.~C.~Myers,
\emph{Towards a derivation of holographic entanglement entropy},
JHEP \textbf{1105}, 036 (2011),
arXiv:1102.0440.

\end{thebibliography}
 
\end{document}
